\documentclass[aspectratio=169]{beamer}
\usepackage[utf8]{inputenc}
\usepackage[T1]{fontenc}
\usepackage[brazil]{babel}
\usepackage{lmodern}
\usepackage{booktabs}
\usepackage{tikz}
\usetikzlibrary{positioning}

\usetheme{Madrid}
\usecolortheme{whale}
\setbeamertemplate{navigation symbols}{}

\title[Empacotamento 2D com Conflitos]{Empacotamento Bidimensional com Conflitos}
\subtitle{Um \textit{solver} guloso e BRKGA para minimizar recipientes}
\author{Ycaro Sales}
\date{\today}

\begin{document}

\frame{\titlepage}

\begin{frame}{Roteiro}
  \tableofcontents
\end{frame}

% ====================================================================
\section{Introdução}

\begin{frame}{A importância da logística no mundo atual}
  \begin{itemize}
    \item A logística é a espinha dorsal da economia moderna: transporte,
          armazenamento e distribuição respondem por boa parte dos custos
          operacionais de empresas.
    \item O \textbf{e-commerce} multiplicou o volume de remessas e a pressão por
          entregas rápidas e baratas.
    \item Acomodar bem as mercadorias --- em caixas, paletes, contêineres e
          caminhões --- é decisivo para a competitividade.
  \end{itemize}
  \vspace{0.5em}
  \begin{block}{Por que minimizar recipientes?}
    Cada contêiner, palete ou caixa economizada significa \textbf{menos
    viagens}, \textbf{menos área de estoque}, \textbf{menos combustível} e
    \textbf{menor impacto ambiental}. Otimizar o empacotamento gera ganho
    econômico e sustentabilidade.
  \end{block}
\end{frame}

% ====================================================================
\section{Definição do problema}

\begin{frame}{Empacotamento Bidimensional (2D Bin Packing) com Conflitos}
  \textbf{Entrada:}
  \begin{itemize}
    \item Recipientes (faixas) retangulares de largura $W$ e altura $H$ fixas;
    \item $n$ itens retangulares, cada um com largura $w_i$, altura $h_i$ e
          classe de cliente $c_i$;
    \item Um \textbf{grafo de conflitos} $G=(\Re, E)$.
  \end{itemize}
  \vspace{0.4em}
  \textbf{Objetivo:} empacotar \emph{todos} os itens minimizando o número de
  recipientes utilizados.
  \vspace{0.4em}
  \textbf{Restrições:}
  \begin{itemize}
    \item Sem sobreposição; item inteiramente dentro do recipiente;
    \item Itens ligados por uma aresta em $G$ \textbf{não podem} dividir o mesmo
          recipiente;
    \item Permite-se rotação de $90^\circ$.
  \end{itemize}
\end{frame}

\begin{frame}{O que são ``conflitos''?}
  \begin{columns}
    \begin{column}{0.55\textwidth}
      Dois itens \emph{em conflito} não podem viajar no mesmo recipiente.
      Exemplos reais:
      \begin{itemize}
        \item Produtos químicos incompatíveis;
        \item Mercadorias de clientes ou rotas diferentes;
        \item Itens que, por segurança/regulação, não podem ser transportados
              juntos.
      \end{itemize}
      \vspace{0.5em}
      O problema é \textbf{NP-difícil}: generaliza o \textit{Bin Packing} e se
      relaciona à coloração de grafos (cada cor $\approx$ um recipiente).
    \end{column}
    \begin{column}{0.45\textwidth}
      \centering
      \begin{tikzpicture}[scale=0.9]
        \node[circle,draw,fill=red!20] (a) at (0,2) {A};
        \node[circle,draw,fill=blue!20] (b) at (2,2) {B};
        \node[circle,draw,fill=green!20] (c) at (1,0) {C};
        \draw[thick] (a) -- (b);
        \draw[thick] (a) -- (c);
        \node at (1,-1) {\small A--B e A--C em conflito};
      \end{tikzpicture}
    \end{column}
  \end{columns}
\end{frame}

% ====================================================================
\section{Representação do problema}

\begin{frame}{Como o problema é modelado}
  \begin{columns}
    \begin{column}{0.56\textwidth}
      Uma \textbf{instância} reúne todos os dados de entrada:
      \begin{itemize}
        \item \textbf{Faixa}: dimensões fixas $W \times H$;
        \item \textbf{Itens}: lista de retângulos, cada um com largura $w_i$,
              altura $h_i$ e classe de cliente $c_i$;
        \item \textbf{Grafo de conflitos} $G$: indexado pela \emph{posição} de
              cada item, com um nó por item e uma aresta por par incompatível.
      \end{itemize}
      \vspace{0.4em}
      As instâncias são lidas de arquivos \texttt{JSON}; a demanda de cada tipo
      de item é expandida em itens individuais.
    \end{column}
    \begin{column}{0.44\textwidth}
      \centering
      \begin{tikzpicture}[scale=0.8]
        \draw[thick] (0,0) rectangle (3,3);
        \node at (1.5,3.3) {\footnotesize faixa $W\times H$};
        \fill[red!25,draw] (0,0) rectangle (1.2,1);
        \node at (0.6,0.5) {\tiny $r_0$};
        \fill[blue!25,draw] (1.2,0) rectangle (2,1.6);
        \node at (1.6,0.8) {\tiny $r_1$};
        \fill[green!25,draw] (0,1) rectangle (1,2.2);
        \node at (0.5,1.6) {\tiny $r_2$};
      \end{tikzpicture}
    \end{column}
  \end{columns}
\end{frame}

\begin{frame}{Representação da solução}
  Uma solução é o conjunto de \textbf{faixas}, cada uma com a posição (origem,
  dimensões e orientação) dos itens nela alocados.
  \begin{block}{Representação geométrica --- \texttt{Container} + EMS}
    \begin{itemize}
      \item Cada faixa mantém seus \textbf{Espaços Máximos Vazios} (retângulos
            livres maximais) e as colocações já feitas;
      \item Inserção pela regra \textbf{bottom-left} (menor origem em $y$, depois
            em $x$), recortando os EMS interceptados;
      \item Checagem de \textbf{conflito}: rejeita o item se um vizinho seu em
            $G$ já está na faixa.
    \end{itemize}
  \end{block}
  \begin{block}{Representação para a busca --- cromossomo BRKGA}
    Vetor de $2n$ \textbf{chaves aleatórias} em $[0,1]$: as primeiras $n$
    codificam a \emph{ordem de inserção} e as últimas $n$ a \emph{orientação} de
    cada item (detalhado adiante).
  \end{block}
\end{frame}

% ====================================================================
\section{Heurísticas utilizadas}

\begin{frame}{Heurística gulosa: First-Fit-Decreasing (linha de base)}
  \begin{enumerate}
    \item Ordena os itens por \textbf{área decrescente};
    \item Insere cada item por \textbf{First-Fit} nas faixas abertas;
    \item Abre uma nova faixa quando nenhuma aberta aceita o item.
  \end{enumerate}
  \vspace{0.5em}
  \begin{block}{Detalhes geométricos}
    \begin{itemize}
      \item Inserção \textbf{ciente de conflitos} (rejeita item se um vizinho em
            $G$ já está na faixa) e com \textbf{auto-rotação}.
      \item Posicionamento via \textbf{Espaços Máximos Vazios} (EMS): o espaço
            livre é um conjunto de retângulos vazios maximais, recortados a cada
            inserção (regra \textit{bottom-left}).
    \end{itemize}
  \end{block}
\end{frame}

\begin{frame}{Metaheurística BRKGA}
  \textbf{Biased Random-Key Genetic Algorithm}: cada solução é um vetor de
  $2n$ chaves aleatórias em $[0,1]$.
  \begin{itemize}
    \item Primeiras $n$ chaves $\rightarrow$ \textbf{ordem de inserção}
          (ordenação crescente);
    \item Últimas $n$ chaves $\rightarrow$ \textbf{orientação}
          ($\ge 0{,}5 \Rightarrow$ rotacionado).
  \end{itemize}
  Um \textbf{decodificador} determinístico empacota por First-Fit honrando os
  conflitos. Aptidão a \emph{minimizar}:
  \[
    f = \texttt{faixas\_usadas} + \texttt{nao\_alocados}\times(n+1)
  \]
  A penalidade $(n+1)$ garante que toda solução inviável seja pior que qualquer
  viável.
  \vspace{0.4em}
  \footnotesize Parâmetros: população $100$; parada com $200$ gerações
  estagnadas; seleção por torneio; \textit{crossover} uniforme; mutação de gene
  único.
\end{frame}

\begin{frame}{BRKGA: decodificação chaves $\rightarrow$ solução}
  \centering
  \footnotesize Exemplo com $n=3$ e chaves $[\,0.9,\ 0.1,\ 0.5\ |\ 0.2,\ 0.5,\ 0.8\,]$:
  \vspace{0.6em}

  \begin{tikzpicture}[node distance=0.4cm]
    \node[draw,fill=blue!10] (ordem) {Ordem (1\textordmasculine\ metade): item1 $<$ item2 $<$ item0 $\Rightarrow$ [1, 2, 0]};
    \node[draw,fill=green!10,below=of ordem] (orient) {Orientação (2\textordmasculine\ metade): item0=normal, item1=rotacionado, item2=rotacionado};
    \node[draw,fill=orange!15,below=of orient] (pack) {First-Fit ciente de conflitos $\Rightarrow$ faixas com itens posicionados};
    \draw[->,thick] (ordem) -- (orient);
    \draw[->,thick] (orient) -- (pack);
  \end{tikzpicture}
  \vspace{0.6em}

  \normalsize Apenas o decodificador conhece o problema; seleção, \textit{crossover}
  e mutação operam genericamente sobre as chaves.
\end{frame}

% ====================================================================
\section{Resultados}

\begin{frame}{Resultados: gulosa vs. BRKGA (18 instâncias)}
  \centering
  \begin{tabular}{lrrr}
    \toprule
    Instância & $n$ & Gulosa & BRKGA \\
    \midrule
    CJCM-E02N20\_Hard\_c0.3           & 20 & 6 & \textbf{4} \\
    CJCM-E00N23\_HardInfeasible\_c0.5 & 23 & 8 & \textbf{6} \\
    CJCM-E03X18\_HardFeasible\_c0.1   & 18 & 3 & \textbf{2} \\
    \midrule
    \textbf{Total (18 instâncias)}    &    & \textbf{91} & \textbf{78} \\
    \bottomrule
  \end{tabular}
  \vspace{0.8em}
  \begin{itemize}
    \item BRKGA: \textbf{melhor em 11}, empate em 7, \textbf{nunca pior}.
    \item Redução média de $\approx \mathbf{14{,}3\%}$ no número de faixas.
    \item Tempos: gulosa $<0{,}01$\,s; BRKGA $\approx 0{,}07$--$0{,}13$\,s por
          instância.
  \end{itemize}
\end{frame}

\begin{frame}{Efeito da densidade de conflitos}
  \begin{itemize}
    \item Os sufixos \texttt{c0.1}, \texttt{c0.3}, \texttt{c0.5} indicam a
          probabilidade de conflito entre pares de itens.
    \item Mais conflitos $\Rightarrow$ menos itens podem dividir um recipiente
          $\Rightarrow$ \textbf{mais faixas necessárias}.
    \item É justamente nos cenários mais restritos que o BRKGA mais se distancia
          da heurística gulosa --- evidenciando o valor da busca metaheurística.
  \end{itemize}
\end{frame}

% ====================================================================
\section*{Conclusão}

\begin{frame}{Conclusão e trabalhos futuros}
  \begin{block}{Conclusão}
    Combinando representação por chaves aleatórias (BRKGA), decodificador
    First-Fit ciente de conflitos e gestão de espaços livres por EMS, obtivemos
    soluções consistentemente melhores que a heurística gulosa, sem nunca
    piorar nenhuma instância.
  \end{block}
  \begin{block}{Trabalhos futuros}
    \begin{itemize}
      \item Incorporar a restrição de \textbf{descarregamento} por classe de
            cliente;
      \item Avaliar sobre um conjunto maior de instâncias;
      \item Ajuste sistemático dos parâmetros da metaheurística.
    \end{itemize}
  \end{block}
\end{frame}

\begin{frame}
  \centering \Huge Obrigado!\\[0.5em]
  \normalsize Perguntas?
\end{frame}

\end{document}
